\scanpage{46}
% archive sheet number: 45

\underline{16 (suite)}

Plus gén, pour tout mono $S\xrightarrow{i}e_R$ dans $R$,
\inserted{($\forall$ cat. $R$ stable par $\varprojlim$ finies ens. \unsure{petit|petite}\deleted{\unsure{es}})}
d'où sous-catégorie $R_{/S}$ de $R$, \deleted{et sous-catégorie} considérons
(\deleted{\unsure{Obj|les}} $\mathrm{topos}\ R$-\unsure{$\mathcal E$}) les $R\xrightarrow{\varphi}\mathcal E$ \inserted{ex. à g} tels que $\varphi(i)$ un isom.
i.e. $\forall x\in\operatorname{Ob}\mathcal E$, $\operatorname{Hom}(x,-)\circ\varphi:R\underset{\text{ex. à g.}}{\longrightarrow}(\mathrm{Ens})$ a la $=$ propriété.

Or les foncteurs ex. à g. \deleted{\unsure{sur}} $R\to\mathrm{Ens}$ s'identifient (avec renv. des flèches)
aux éléments de $\operatorname{Pro}R$ (sont pro-représentables), et on prend donc les $\xi\in\operatorname{Pro}R$ tels que
$\operatorname{Hom}_{\operatorname{Pro}R}(\xi,S)\ne\varnothing$, i.e. qui sont sur $S$,
donc ce sont ess. les éléments de $\operatorname{Pro}(R_{/S})$, \deleted{donc} ou encore la sous-cat. pleine de
$\mathcal T=(\operatorname{Pro}(R))^\circ=\operatorname{Ind}(\unsure{\mathcal C|\mathcal G|\mathcal S})$ ($\mathcal C=R^\circ$)
\inserted{des $\xi$} qui \deleted{sont} tels qu'il existe un morphisme $\Lambda\to\xi$, où $\Lambda$ est l'objet de $\mathcal C$ correspondant : $S$.

On a pour les structures envisagées le topos \deleted{\unsure{ouvert}}\inserted{classifiant}
$\widehat R_{/S}$, avec $\widehat R_{/S}\xrightarrow{\widehat i}\widehat R$ déduit de $i:R_{/S}\to R$.

\[
\begin{aligned}
 &\text{On a}\quad i_{\unsure{!|1}}^{\deleted{\unsure{*}}}=\text{inclusion de $R_{/S}$ ds $R$},
 \quad i^*(X)=X\times_e S/S\\
 &i_*(\underset{\text{Objet de $R$ sur $S$}}Y)
 =\{ \underset{\operatorname{Ob}R}{X}\longmapsto\operatorname{Hom}_S(X\times S,Y)=\operatorname{Hom}(X\times S,Y)\},\\
 &\text{ou}\quad i_*(\deleted{\unsure{F|Y}}Y)=\underline{\operatorname{Hom}}(S,Y).
 \quad\inserted{\text{variable}}
\end{aligned}
\]

Considérons maintenant une famille $S_\alpha\lhook\joinrel\xrightarrow{i_\alpha}e_R$ \inserted{filtrante} décroissante
de ss-objets de $e_R$, et soit $T'=\bigcap_\alpha T_{S_\alpha}$ \deleted{\unsure{l'espe}} sous-espèce de structure de $T=T_R=T_{\widehat R}$
définie par les conditions $\varphi(i_\alpha)(=f^*(i_\alpha))$, où $\underline{\underline f}:\mathcal E\to B_T=\widehat R$ est isom pour tt $\alpha$.

\uline{Définie} ainsi par la condition sur les $\operatorname{Hom}(\unsure{X|Z|X_\xi},-)\circ\varphi$, on est ramené à examiner le cas du topos ponctuel $\mathcal E=(\mathrm{Ens})$.
\deleted{\unsure{On doit donc considérer les|Oui}} d'où $\xi=(Y_n)\in\operatorname{Ind}\mathcal C$,
$S=(S_\alpha)\in\operatorname{Pro} R$, $\xi\to e_R=e_{\operatorname{Pro} R}$ est un mono ;
les objets envisagés de \unsure{$\mathcal C^\circ$|$\mathcal S$} sont ceux pour les quels il existe un (unique) morphisme $\xi\to\unsure{O|\sigma|0|\mathcal O}$, donc dans
$\operatorname{Pro}R=\mathcal R$ ce sont ceux $x$ qui se factorisent par un (unique)
$x\longrightarrow S\hookrightarrow e_R=e_{\mathcal R}$.
Mais \deleted{\unsure{quelle|quelles}} est la catégorie ens. petite $R'$ avec \uline{lim} classifiante ?
\deleted{\unsure{Le sous-topos|C'est le topos}}\inserted{pour $\widehat{R'}$ est le} $\bigcap_\alpha B_{/S_\alpha}$.
On a déterminé ces pts, mais quels sont ces pts \unsure{ens} ?

% margin note written sideways
\margin{sauf erreur, ce sont les éléments de l'image essentielle de $R\to\operatorname{Pro}R_{/S}$
définis par chgt de base $X\to X\times_e S$, ou encore les éléments de la catégorie
$\varinjlim R_{/S_\alpha}$. On trouve donc
$\bigcap_\alpha B_{T_\alpha}=\bigcap_\alpha\widehat{R_{/S_\alpha}}=\bigcap_\alpha\widehat R_{/S_\alpha}
=\widehat{\varinjlim_\alpha R_{/S_\alpha}}$, la \inserted{petite} catégorie modulaire avec $\varprojlim$
universelle est $\varinjlim R_{/S_\alpha}$ (qui est une catégorie de fractions de $R$ pour les flèches $i$ telles que $i\times_e S$ un iso.
$\widehat{R_{/S_\alpha}}$, $\widehat R_{/S_\alpha}$ plongés dans $B$ par $i_{\alpha *}$, \unsure{non par|un}
\unsure{$i_{\alpha !}!$|$i_{\alpha !}$ !!}}
